All SPSS Modeler streams, datasets, intermediate outputs, and generated
models used in this project are provided in the accompanying
\texttt{SPSS\_modeler} directory. The main SPSS Modeler project file is:

\begin{verbatim}
SPSS_modeler/dugn944.cpj
\end{verbatim}

The project was developed using FlexIT and macOS. Some SPSS Modeler source
nodes retain absolute file paths from the development environments. Therefore,
if a file cannot be located when the project is opened on another computer,
the source node should be redirected to the corresponding file in the
\texttt{SPSS\_modeler} directory. No external dataset is required.

\textbf{Important:} SPSS Modeler may retain absolute paths from the original
development computer. If a source-node path is unavailable, please select the
file with the same filename from the submitted \texttt{SPSS\_modeler}
directory. All required source data are included locally.

\subsection{Required Data Files}

The original dataset is retained unchanged as:

\begin{verbatim}
SPSS_modeler/dataset_gallstone.xlsx
\end{verbatim}

For the data-integration stage, the original dataset was divided into two
logically separate sources:

\begin{verbatim}
SPSS_modeler/gallstone_demographic_clinical.csv
SPSS_modeler/gallstone_biological_body_composition.csv
\end{verbatim}

Both integration files contain 319 participant records and share the same
\texttt{record\_id} values. They are merged using \texttt{record\_id} in the
Data Cleaning/Data Processing workflow.

If SPSS Modeler reports that a source file cannot be found, open the
corresponding source node, choose \textit{Browse}, and select the appropriate
file from the local \texttt{SPSS\_modeler} folder rather than using the
absolute path stored on the original development machine.

\subsection{SPSS Modeler Streams}

The main streams used in the project are:

\begin{verbatim}
Business_Understanding.str
EDA_stream.str
Data_Cleaning.str
Data_Processing.str
Feature_Transformation.str
Feature_Importance.str
Modeling.str
Auto_Modeling.str
\end{verbatim}

Their intended execution order is approximately:

\begin{center}
\begin{tabular}{c}
Business Understanding $\longrightarrow$ EDA
$\longrightarrow$ Data Cleaning and Integration \\
$\longrightarrow$ Transformation $\longrightarrow$ Feature Selection \\
$\longrightarrow$ Modelling $\longrightarrow$ Auto Modelling
\end{tabular}
\end{center}

The streams ending in \texttt{.str-} are backup versions automatically
created during development and are not required for normal execution.

\subsection{Generated and Intermediate Outputs}

The directory also contains generated datasets, models, and analytical
outputs used to support the report. Important files include:

\begin{verbatim}
Output_Feature_Selection.tab
Transformed_Data.tab
GeneratedModels.gen
new.gen
Correlation matrix.cou
Data Audit of [39 fields].cou
t-SNE.cou
Transformed_Data (43 fields, 319 records).cou
\end{verbatim}

Files with the \texttt{.cou} extension contain SPSS Modeler output objects,
while \texttt{.gen} files contain generated model objects. These files are
included so that the modelling outputs can be inspected without necessarily
rebuilding every model from the beginning.

\subsection{Reproducing the Workflow}

To reproduce the analysis:

\begin{enumerate}
    \item Open \texttt{dugn944.cpj} in SPSS Modeler.
    \item Confirm that the source nodes point to files inside the submitted
    \texttt{SPSS\_modeler} directory.
    \item If an absolute-path error occurs, manually reselect
    \texttt{dataset\_gallstone.xlsx} or the appropriate integration CSV file
    from this directory.
    \item Run the relevant stream from left to right.
    \item For the integration workflow, verify that
    \texttt{record\_id} is used as the merge key.
    \item Confirm that the integrated dataset contains 319 records and that
    no additional missing values or duplicate records have been introduced.
    \item Run the feature-transformation and feature-selection streams before
    executing the refined modelling workflows where required.
    \item Use the existing generated model and output files to compare the
    reproduced results with those reported in the main document.
\end{enumerate}

The modelling dataset is partitioned into training, validation, and testing
subsets within SPSS Modeler. Therefore, the submitted files should be used
with the existing partition settings and random seed where specified to
maintain consistency with the results reported in this project.

\subsection{Recommended Submission Directory}

The submitted project is organised as follows:

\begin{verbatim}
submission_directory/
|
|-- SPSS_modeler/
|   |
|   |-- dugn944.cpj
|
|   |-- dataset_gallstone.xlsx
|   |-- gallstone_demographic_clinical.csv
|   |-- gallstone_biological_body_composition.csv
|
|   |-- Business_Understanding.str
|   |-- EDA_stream.str
|   |-- Data_Cleaning.str
|   |-- Data_Processing.str
|   |-- Feature_Transformation.str
|   |-- Feature_Importance.str
|   |-- Modeling.str
|   |-- Auto_Modeling.str
|
|   |-- Output_Feature_Selection.tab
|   |-- Transformed_Data.tab
|   |-- GeneratedModels.gen
|   |-- new.gen
|
|   |-- *.cou
|   `-- split_gallstone_dataset.py
|
|-- Report/
|   |-- main.tex
|   |-- references.bib
|   |-- media/
|   `-- other report source files
\end{verbatim}

Keeping all project resources under a single parent directory reduces the
risk of missing dependencies and allows the marker to redirect any unresolved
absolute paths to the equivalent submitted files.
